% 精读章: Hsiang, Climate Econometrics (Annual Review of Resource Economics 2016)
% 原文出处: doi:10.1146/annurev-resource-100815-095343（所据版本：NBER Working
%          Paper 22181, April 2016——本仓 papers/01_methods/2016-NBER-w22181-*.pdf
%          为团队正典，本章页码引注均指该 NBER 版的印刷页码；已刊版 ARRE 8: 43--75
%          并存，页码不同，引用页码须注明所据版本。期刊卷期与 DOI 不见于 NBER 版
%          PDF，出处见 bibs/refs.bib 该条的 note）
% 体例注记: 2026-10-21 起，正文中的 terms/expressions 一律用英文原词，
%          中文仅保留句法成分；定义处随文给中文对照并登记英文 \index。
% 建档/更新日期由章首 \chdate 渲染（grep: chdate）
% 主笔: 待认领；起草: Claude（2026-11-04，bjr 审订前稿）
% 核对: 2026-11-19 对照本地 PDF 全章重核——原句引用逐条改回原文字面；
%       改正式 (16)--(20) 的推导次序、式 (24)、玉米应用的频段、骨架表各节
%       页码，以及 marginal treatment comparability 与 unit homogeneity 的关系。
% partinfo 定位: 把「weather 回归能否识别 climate 效应」的争论正式化，整理
%               cross-section／time-series／long differences 三种 research designs
%               与两件工具（marginal treatment comparability、
%               frequency-identification tradeoff），见
%               \chref{ch:info-climate-change-adaptation} 的方法线与速查表。

\chapter{Climate Econometrics（Hsiang，ARRE 2016）}
\label{ch:hsiang-2016}
\chdate{2026-11-04}{2026-11-19}

\section{这篇论文在说什么}

一句话概括：这篇综述（\citealp{hsiang2016climate}，所据 NBER w22181 版）把``regressions on `weather' variables''\citep[第~2 页]{hsiang2016climate}能否识别 climate 效应这个``particularly thorny''\citep[第~2 页]{hsiang2016climate}的争论正式化，在可以写明的条件下得出``the use of weather variables in regressions is justified and, perhaps surprisingly, dominates traditionally preferred methods''\citep[第~2 页]{hsiang2016climate}；同时整理了 climate econometrics 的三种 research designs 与两件工具（frequency-identification tradeoff、对 marginal treatment comparability 的部分检验），并给出 weather variation 恰好识别 climate 效应的充分条件（Envelope Theorem 与 Gradient Theorem，\S 2.4）。

它在 partinfo 里的坐标：方法线上联结两条线的一站（\chref{ch:info-climate-change-adaptation}）——它把 cross-sectional line 与 panel line 的设计放进同一组记号，逐一列出各自依赖的 assumption。本讲义的归纳：\chref{ch:deschenes-greenstone-2007} 自认的局限与 \chref{ch:burke-hsiang-miguel-2015} 的 adaptation 判断，都可以在这套框架里重新表述；\chref{ch:merel-gammans-2021} 随后自陈与其关系：``our result is both different and more specific''\citep[第~1217 页]{merel2021climate}。

骨架表（所据版本：NBER Working Paper 22181；页码为该版印刷页码。\S 4 及以后本讲义只读到与本 part 直接相关的 \S 4.1.3，其余待补读）：

\begin{table}[htbp]
\centering\small
\begin{tabular}{@{}p{4.4cm}p{1.4cm}p{7.2cm}@{}}
\toprule
原文章节 & 页码 & 讲什么 \\
\midrule
\S 1 Introduction & 2--4 & 1.1 climate 的形式化定义（$\mathbf{C}_{i\tau}$ 对 $\mathbf{c}_{i\tau}$）；1.2 direct 与 belief 两条作用通道 \\
\S 2 The Empirical Problem & 4--14 & 2.1 三种 research designs；2.2 低频 variation 与可信识别的权衡；2.3 marginal treatment comparability 的部分检验；2.4 用 weather variation 恰好识别 climate 效应 \\
\S 3 Measurement of Climate Variables & 15--17 & 从社会与经济的角度压缩 climate 变量的维度；测度创新（Figure~4） \\
\S 4 Econometric Models & 17--25 & 4.1 非线性效应；4.2 位移与滞后；4.3 统计不确定性；4.4 adaptation 的测度；4.5 跨研究的比较与综合 \\
\S 5 Attribution and Projection & 25--26 & 历史归因；未来 climate change 效应的 projection \\
\S 6 Remaining Challenges & 26--27 & 效应与机制的对应、adaptation 与一般均衡、无先例的事件、与理论和数值模型的结合 \\
\S 7 Conclusion & 27--28 & 结论 \\
\bottomrule
\end{tabular}
\end{table}

\section{预备知识}

三件工具，首用交代来历。其一，因果推断的 Fundamental Problem of Causal Inference（Holland 1986，因果推断文献的术语）：同一总体不可能在同一时段同时经历 $\mathbf{C}$ 与 $\mathbf{C}+\Delta\mathbf{C}$，故式~\ref{eq:hs-ate} 的处理效应不可直接观测（\S 2 开头，第~4 页）。其二，\index{envelope theorem}Envelope Theorem（包络定理，微观经济学的标准结果）：最优值函数对参数的导数，等于目标函数在最优点处对该参数的偏导数；选择变量随参数的调整对最优值没有一阶影响。其三，\index{gradient theorem}Gradient Theorem（梯度定理，多元微积分的标准结果）：梯度沿一条路径的线积分等于路径两端函数值之差。后两者在 \S 2.4 使用（第~10--14 页；论文在第~11 页点出这两个定理的名字，但不作陈述，以上表述是教科书的通行说法）。

\section{数据与测度}

综述不依赖单一数据集，本节整理它对「climate 变量怎么测」的要点（\S 3，第~15--17 页）。\S 3 的主题是如何从社会与经济的角度压缩 $\mathbf{c}$ 的维度。早期研究多用粗测度（纬度、事件虚拟变量），作者列举其缺陷，其中之一是``coarse climate measures introduce large measurement errors that will cause attenuation bias, leading to under-rejection of the null hypothesis''\citep[第~15 页]{hsiang2016climate}；另一条是研究者常常关心 \index{dose-response function}dose-response function（剂量—反应函数）$\E[Y\mid\mathbf{c}] = f(\mathbf{c})$ 的结构（式 (25)，第~15 页），粗测度给不出。测度创新的方向见其 Figure~4（第~15 页图注）：用逐日最低、最高温度之间插值出的逐小时温度构造 degree-days（Schlenker 与 Roberts 2009）、用风场模型重建热带气旋暴露（Hsiang 与 Jina 2014）、识别与 ENSO 强耦合的 teleconnected 像素（Hsiang 等 2011）、以下雨天数刻画年内降水的分布（Fishman 2016）。这一节与本讲义对读者的要求一致：climate 变量的口径是跨研究不可比的首要来源。

\section{模型设定与识别策略：weather 回归何时识别 climate 效应}

本节是论文的理论贡献，按原文四步走：定义、问题、三种设计的取舍、两件工具与一个正式结果。

\subsection{第一步：climate 与 weather 的形式化定义}

位置 $i$ 在每个时刻 $t$ 有一个状态向量 $\mathbf{v}_{it} = [$temperature, precipitation, humidity,\dots$]$（式 (1)，第~2 页）；在时段 $\tau$ 内，它从分布 $\psi(\mathbf{C}_{i\tau})$ 中抽出（式 (2)）。刻画该分布的参数向量 $\mathbf{C}_{i\tau}$ 定义为 \index{climate}climate；时段内实际实现的状态有一个经验分布 $\psi(\mathbf{c}_{i\tau})$，刻画它的向量 $\mathbf{c}_{i\tau}$ 定义为 \index{weather}weather（第~2--3 页）——``Define $\mathbf{C}_{i\tau}$ to be the \emph{climate} at $i$ during $\tau$, since it characterizes the distribution of possible realized states $\mathbf{v}_{it}$''\citep[第~2 页]{hsiang2016climate}，``we define $\mathbf{c}_{i\tau}$ to be a description of the \emph{weather} during $\tau$''\citep[第~3 页]{hsiang2016climate}。本 part 的 weather／climate 区分（\chref{ch:info-climate-change-adaptation}）在计量上对应的就是这一对记号。作者还强调 climate 没有固有的时间尺度（通行的 30 年平均定义``is fairly arbitrary''\citep[第~3 页]{hsiang2016climate}）。

climate 影响结果有两条通道：\index{direct effect}direct effect 与 \index{belief effect}belief effect。``the climate during $\tau$ influences what realizations of weather $\mathbf{c}$ actually occur during that interval, which in turn affects a population directly (e.g. a rainy climate generates rain, causing people to get wet)—call this the `direct effect' of climate. Second, individuals' beliefs over the structure of $\mathbf{C}$ may affect their decisions and resulting outcomes, regardless of what $\mathbf{c}$ is realized (e.g. if people believe their climate is rainy, some will buy umbrellas)—refer to this as the `belief effect'''\citep[第~3 页]{hsiang2016climate}。总效应写成 $Y(\mathbf{C}) = Y(\mathbf{c}(\mathbf{C}), \mathbf{b}(\mathbf{C}))$（式 (3)，第~3 页），其梯度按 direct 与 belief 两块分解（式 (4)，第~4 页）；作者注：``interactions between beliefs and direct impacts ($\frac{\partial^2 Y}{\partial b_n\partial c_k}$) and belief effects themselves are together often referred to as `adaptations' in the literature''\citep[第~4 页]{hsiang2016climate}——在这个框架里，adaptation 就是 belief effect 加上它与 direct effect 的交互。

\subsection{第二步：实证问题与两条 assumption}

目标是 climate 变化的平均处理效应（式 (5)，第~4 页）
\begin{equation}
\beta = \E[Y_{i\tau}\mid\mathbf{C}_{i\tau}+\Delta\mathbf{C}_{i\tau}, \mathbf{x}_{i\tau}] - \E[Y_{i\tau}\mid\mathbf{C}_{i\tau}, \mathbf{x}_{i\tau}] ,
\label{eq:hs-ate}
\end{equation}
其中 $\mathbf{x}$ 是可观测的非 climate 因素。同一单元不能同时处于两种 climate 之下，只能拿另一个单元或另一个时段来近似。跨单元比较要求两单元同质：$\E[Y_{i\tau}\mid\mathbf{C},\mathbf{x}_{i\tau}] = \E[Y_{j\tau}\mid\mathbf{C},\mathbf{x}_{j\tau}]$，即 \index{unit homogeneity}unit homogeneity assumption（单位同质假设，式 (6)，第~4--5 页）。时间序列路线只需要它的一个较弱的形式（同一单元跨时期可比），但必须另加第二条 assumption：\index{marginal treatment comparability}marginal treatment comparability assumption（边际处理可比性，第~6 页）——某一时段实现的 weather 相对基准 climate $\mathbf{C}_1$ 的偏离 $\mathbf{c}_{\tau}-\mathbf{C}_1$ 所带来的结果变化，等于 climate 参数本身作同样大小的调整 $\Delta\mathbf{C} = \mathbf{C}_2-\mathbf{C}_1$ 所带来的结果变化（式 (9)，第~6 页）。作者直言这条 assumption``has been widely debated''\citep[第~6 页]{hsiang2016climate}，故后文给检验与条件。

\subsection{第三步：三种 research designs 的取舍}

``There are essentially three research designs in use that approximate the average treatment effect in Eq.~5: cross-sectional approaches, use of time-series variation, and a hybrid known as `long differences'''\citep[第~5 页]{hsiang2016climate}。各自的 variation 与代价：

\textbf{\index{cross-sectional approach}Cross-sectional approach。}同一时期跨单元比较（式 (8)）。长处是隐含地包括了 adaptation：``since farmers who inhabit a location for a long period will have a strong grasp of $\mathbf{C}$ at their location and will adjust farm investments and management to optimize based on these beliefs, farm prices can be assumed to reflect all direct effects as well as all belief effects''\citep[第~5 页；此处在原文中即对 Mendelsohn 等 1994 农地价值路线的概括]{hsiang2016climate}。短处：``A weakness of the cross-sectional approach is its vulnerability to omitted variables bias''\citep[第~5 页]{hsiang2016climate}，且``there exists no systematic method for determining whether any key variables are omitted from Eq.~8 and thus an econometrician can never be certain their model is unbiased''\citep[第~5 页]{hsiang2016climate}。

\textbf{Time-series（含 panel）approach。}单元自身跨期比较（式 (10)，unit fixed effects 加趋势项）。这一思路由 Huntington（1922）提出，他还用工厂工人的数据实施过类似的比较（第~6--7 页）；它获得现代地位是在 Deschênes 与 Greenstone（2007）手中：``This approach gained prominence in modern economic analysis when used by Deschênes \& Greenstone (2007), who analyzed whether agricultural profits in US counties responded to `random fluctuations in weather'''\citep[第~7 页]{hsiang2016climate}。长处：``The core benefit of this approach is that it accounts for unobservable differences between units, eliminating a potential source of omitted variables bias''\citep[第~7 页]{hsiang2016climate}。短处：随时间变化、又与 $\mathbf{c}$ 相关的因素若被遗漏，估计仍会有偏；作者引 Hsiang（2010）的例子——温度、降水与飓风都受 ENSO 调制，漏掉其中重要的 climate 变量，$\hat\beta_{TS}$ 就可能有偏（第~7 页）。

\textbf{\index{long differences}Long differences。}跨单元比较各自在长时段上的变化（式 (11)，第~8 页），定位是折中：``An approach that aims to compromise between the strengths and weaknesses of cross-sectional analysis and time series identification is the long differences strategy, where changes for both the outcome and the climate within locations are correlated across locations''\citep[第~7 页]{hsiang2016climate}；在增长、作物单产与冲突的应用中，研究者``have found that $\beta_{LD}$ is almost identical to $\beta_{TS}$, leading them to conclude that gradual changes in $\mathbf{c}$ likely induce similar effects to more rapid changes in $\mathbf{c}$''\citep[第~8 页，指 Dell et al. 2012、Burke \& Emerick、Lobell \& Asner 2003、Burke et al. 2015b 等的应用]{hsiang2016climate}。

\subsection{第四步：两件工具与一个正式结果}

\textbf{工具一：\index{frequency-identification tradeoff}frequency-identification tradeoff。}belief effect 只对持续的环境变化起作用——只持续五分钟的变化不会让人更新对 climate 的信念（第~8 页）——所以``it is widely thought that low-frequency data ... is required to measure belief effects when using time-series variation''\citep[第~8 页]{hsiang2016climate}；但所用 variation 的频率越低，unit homogeneity 越难成立（相隔数世纪的同一总体未必可比），``it comes at the cost of less credible identification, an issue highlighted by Hsiang \& Burke (2014) as the \emph{frequency-identification tradeoff}''\citep[第~8--9 页]{hsiang2016climate}。缓解条件：若 belief effects 主要通过改变 direct effects 起作用，``total effects of climate still may be nearly identified with high-frequency time-series''\citep[第~9 页]{hsiang2016climate}。

\textbf{工具二：部分检验（\S 2.3）。}若三种设计的估计相等，$\hat\beta_{CS}=\hat\beta_{LD}=\hat\beta_{TS}$，则``one possible explanation is that the marginal treatment comparability assumption is valid and temporary changes in realizations of $\mathbf{c}$ have similar effects to analogous changes in $\mathbf{C}$''\citep[第~9 页]{hsiang2016climate}。它的实施见下文「结果与解释」一节。

\textbf{正式结果：exact identification（\S 2.4）。}这是本综述最重的理论贡献。设结果是最大化问题的值函数（式 (14)，第~11 页）
\begin{equation*}
Y(\mathbf{C}) = \max_{\mathbf{b}\in\mathbb{R}^N} z\bigl(\mathbf{b}, \mathbf{c}(\mathbf{C})\bigr),
\end{equation*}
$z$ 对 $\mathbf{b}$ 可微且为凹。推导分五步（式 (15)--(20)，第~11、13 页）。先对 climate 的第 $k$ 个分量用链式法则，$\mathrm{d}Y/\mathrm{d}C_k$ 拆成三项：$z$ 对 $C_k$ 的直接偏导、经由行动 $b_n$ 的间接项、经由 weather $c_\kappa$ 的间接项（式 (15)）。第一项为零：climate 作为一个概率分布，只能经由它引起的 weather 实现和基于信念的行动影响结果，$\partial z/\partial C_k=0$（式 (16)）。第二项为零：$Y$ 是 $z$ 对全部 adaptation 行动优化后的结果，最优处的一阶条件给出 $\partial z/\partial b_n=0$（式 (17)）。代入后只剩经由 weather 的一项（式 (18)，作者称之为 Envelope Theorem 的应用）。最后只看彼此对应的 weather 与 climate 维度，$\mathrm{d}c_\kappa/\mathrm{d}C_k$ 在 $\kappa=k$ 时为 1、否则为 0（式 (19)），得到式 (20)：
\begin{equation*}
\frac{\mathrm{d}Y(\mathbf{C}_1)}{\mathrm{d}\mathbf{C}_k} = \frac{\partial Y(\mathbf{C}_1)}{\partial \mathbf{c}_k} ,
\end{equation*}
即``Locally, the marginal effect of the climate on $Y$ is identical to the marginal effect of the weather. Eq.~20 implies that Eq.~9, the marginal treatment comparability assumption, holds''\citep[式 (20) 及引文，第~13 页]{hsiang2016climate}。再用 Gradient Theorem 把 marginal effect（边际效应）沿 climate 路径积分，得非边际 climate 变化的效应（式 (23)，第~14 页）
\begin{equation}
Y(\mathbf{C}_2) - Y(\mathbf{C}_1) = \int_{\mathbf{C}_1}^{\mathbf{C}_2} \hat\beta_{TS}\big|_{\mathbf{C}} \cdot \mathrm{d}\mathbf{C} ,
\label{eq:hs-integral}
\end{equation}
``in the presence of adaptation adjustments in $\mathbf{b}$, using only time-series estimates''\citep[式 (23) 及引文，第~14 页]{hsiang2016climate}。作者的总结：``If the outcome is a solution to a maximization problem (Eq.~14) for a function $z(.)$ that is continuous and differentiable in the space of all adaptive actions $\mathbf{b}$, then by application of the Envelope Theorem (Eq.~18) we know that the marginal effect of the climate is exactly the same as the marginal effect of an equally structured change in the weather distribution (Eq.~20), if both are evaluated locally relative to an initial climate. By the Gradient Theorem we know that a sequence of marginal effects of the weather empirically estimated via time-series variation at sequential values of $\mathbf{C}$ can then be integrated to compute the effect of non-marginal climate changes''\citep[第~14 页]{hsiang2016climate}。

这一结果的直觉是论文自带的相邻两户例子（第~11 页；几何版本是其 Figure~3，第~12 页）：一条南北向街道上的两户邻居，其他方面完全相同，只是纬度略异、climate 略异，各自都已针对自己的 climate 做了优化。两户 adaptation 决策的细微差别带来的净效应为零，因为北户多得的边际收益恰被多付的边际成本抵消（两户都处在最大值点）。于是南户若经历北户的 weather，其 weather 效应就等于两户之间的截面差异。\textbf{注意适用范围}：这一等价只在「结果已被优化」的边际处成立；截面路线要恢复 $Y(\mathbf{C})$，还需各单元的积分常数相同（式 (24)，$\mathrm{d}\phi/\mathrm{d}i=0$），``which implies the strong form of the unit homogeneity assumption that units are comparable in levels conditional on the climate (Eq.~6). Thus the set of assumptions necessary for valid cross-sectional identification in this setting is strictly larger than the set of assumptions required for valid time-series identification''\citep[第~14 页]{hsiang2016climate}。

\begin{intuition}[Envelope 论证与 M-G 切点的关系]
Envelope 论证说的是：在一个已经优化的 climate 处，weather 的 marginal effect 等于 climate 的 marginal effect。对「一条在 weather 实现上非线性的面板回归」能恢复什么，综述在 \S 4.1.3 另有讨论（第~19--20 页）：若存在一条全局的非线性响应，这样的回归可以把它估出来，但响应面只在各单元 weather 实现的切点附近由时间序列 variation 局部识别，结论在 weather 实现趋于本地均值的极限下才严格成立。\chref{ch:merel-gammans-2021} 的几何把这一限定讲得更具体：在他们的 climate penalty 模型里，各 location 的 short-run response 在自己的 $\mu_i$ 处与 long-run response 相切（该章图~\ref{fig:cca-lrsr}），切点就是 envelope 条件成立的地方；离开切点，漏掉 penalty term 的面板系数是两种 response 的加权平均（自陈：``our result is both different and more specific''\citep[第~1217 页]{merel2021climate}）。两文合读，这一论证的适用边界才完整。
\end{intuition}

\section{结果与解释：部分检验的实施（玉米单产）}

综述自己的经验证据集中在 \S 2.3 的应用（第~10 页；Figure~2 在第~9 页）。数据：Schlenker 与 Roberts（2009）所用的县级玉米单产、温度与降水数据，更新至 2014 年，限于 100 度经线以东、没有缺失观测的 730 个县。做法：对三个变量的序列施 Baxter--King 带通滤波，按 2--5、6--9、10--13、14--17、18--33 年五个周期段分别估计式 (10)（滤波样本为 1962--2002 年），与原始年度数据的估计（$\hat\beta_{TS}$，即 Schlenker 与 Roberts 的做法）、long differences（$\hat\beta_{LD}$，Burke 与 Emerick）及横截面（$\hat\beta_{CS}$，Schlenker 等 2006）并列。结果：``In all cases, except the cross-section, these estimated effects are near one another and not statistically different, suggesting that variations in temperature over time have similar effects on maize yields in this context, regardless of the time-scale of these variations''\citep[第~10 页]{hsiang2016climate}。横截面例外的原因可辨：``the fact that $\hat\beta_{CS}$ changes substantially (to more closely resemble time-series estimates) when rainfall terms are \emph{omitted} highlights the vulnerability of the cross-sectional approach to misspecification''\citep[第~10 页]{hsiang2016climate}。作者小结：``Nonetheless, these results overall appear consistent with an assumption of marginal treatment comparability in this context, at least for time scales shorter than 33 years''\citep[第~10 页]{hsiang2016climate}。\textbf{读图须知}：检验的说服力直接受各设计的 unit homogeneity assumption 是否成立的约束，``It is theoretically possible to obtain $\hat\beta_{CS}=\hat\beta_{LD}=\hat\beta_{TS}$ by chance even if all key assumptions are violated''\citep[第~10 页]{hsiang2016climate}——各估计相等不能证明 assumption 成立，只能排除一部分反驳。

\section{局限与批评}

论文自己承认的：\S 2.4 的结论依赖三个条件，作者指出既有的经验研究并不检查它们——``that $Y$ is the solution to a [constrained] maximization, that adaptations $\mathbf{b}$ take on continuous values, and that the maximand function $z(.)$ is differentiable in $\mathbf{b}$''\citep[第~14 页]{hsiang2016climate}；而且``many prior studies do not properly compute climate effects via Eq.~23''\citep[第~14 页，作者举 Schlenker et al. 2013 与 Houser et al. 2015 为突出的例外，并提到另有数项研究计算正确或以线性外推近似]{hsiang2016climate}；部分检验中各估计相等可以出于巧合（第~10 页）。\S 6 登记的其余方法论挑战（第~26--27 页）本讲义未读，待补写。后续文献：\chref{ch:merel-gammans-2021} 指出边际处等同只是特例——漏掉 penalty term 的二次面板估计到的是 long-run 与 short-run response 的加权平均，与 long-run response 相切的位置由 weather 的 variation 结构决定（该文第~1215--1219 页）。

本讲义的评注：本综述与 \chref{ch:merel-gammans-2021} 是一组必须对着读的文献。Hsiang 回答「何时可以相信 weather 回归识别了 climate 效应」（边际处、被优化的结果），M-G 回答「在一个具体的函数形式下，偏离有多大、偏向哪边」（由 skewness 与 variation 构成决定）。前者的 Envelope 论证对应后者的切点几何；后者的偏差公式是前者没有给出的量化。

\section{在本领域脉络中的位置}

本讲义的归纳：这篇综述上承两站经验研究留下的问题——\chref{ch:deschenes-greenstone-2007} 自认面板估计不含全部 long-run 调整，\chref{ch:burke-hsiang-miguel-2015} 的非线性面板系数该如何解读尚无定论——它把这些问题放进同一框架，给出三种 designs 各自的 assumption 清单、一个部分检验与一个正式结果。它开的新问题：frequency-identification tradeoff 之下，adaptation 的直接测度如何做（原文 \S 4.4，第~23--24 页，待补读）；\chref{ch:merel-gammans-2021} 的偏差分解可视为在二次设定下对「面板系数含多少 long-run 成分」的第一个闭式回答。本章在 \chref{ch:info-climate-change-adaptation} 的方法线与速查表中均有登记。

% 章末参考文献（只列本章引用的条目），勿删
\chapterrefs
